\chapter{学習順序を頻度と転用性から組み立てる}
大学院物理の背景がある読者は、既知の力学・熱・電磁気を原典の条件確認で回収し、情報・化学・地球科学・制度の新規概念へ時間を配分する。次の列挙は出題数による着手候補であり、習得時間を測ったS/A/B/Cの確定優先度ではない。難易度は校準初稿、学習時間と得点効率は未評価として残す。
\section{基礎の頻度上位に接続する基礎}
エネルギー（29問）、生物・バイオ（28問）、科学技術史・社会（26問）、材料・金属（26問）、環境・持続可能性（26問）を最初の点検対象とする。
\section{適性の頻度上位に接続する基礎}
安全・リスク（59問）、コンプライアンス・説明責任（53問）、技術者倫理・意思決定（41問）、技術士法・義務責務（29問）、研究・情報倫理（29問）を最初の点検対象とする。
\section{応用理学の頻度上位に接続する基礎}
物理探査（68問）、物理化学（48問）、電磁気・物性（40問）、有機・高分子化学（33問）、応用地質・地盤（27問）を最初の点検対象とする。
基礎は保存則、単位、論理、確率と待ち行列の成立条件を共通の道具にする。適性は技術士法の義務責務と、公益・安全・説明責任を別の根拠で整理する。応用理学は波の伝播・場・収支を物理探査へつなぎ、鉱物・岩石・地層・年代・同位体を用語と観測証拠から学ぶ。
頻度の低い問題も索引と解説から除外しない。高リスクの法令や歴史統計は解説の留保を先に読み、出題時点の原典と現在の参考説明を分ける。頻度だけで誤った初稿を反復学習しない。
